%% Generated by Sphinx.
\def\sphinxdocclass{report}
\documentclass[letterpaper,10pt,english]{sphinxmanual}
\ifdefined\pdfpxdimen
   \let\sphinxpxdimen\pdfpxdimen\else\newdimen\sphinxpxdimen
\fi \sphinxpxdimen=.75bp\relax
\ifdefined\pdfimageresolution
    \pdfimageresolution= \numexpr \dimexpr1in\relax/\sphinxpxdimen\relax
\fi
%% let collapsible pdf bookmarks panel have high depth per default
\PassOptionsToPackage{bookmarksdepth=5}{hyperref}

\PassOptionsToPackage{booktabs}{sphinx}
\PassOptionsToPackage{colorrows}{sphinx}

\PassOptionsToPackage{warn}{textcomp}
\usepackage[utf8]{inputenc}
\ifdefined\DeclareUnicodeCharacter
% support both utf8 and utf8x syntaxes
  \ifdefined\DeclareUnicodeCharacterAsOptional
    \def\sphinxDUC#1{\DeclareUnicodeCharacter{"#1}}
  \else
    \let\sphinxDUC\DeclareUnicodeCharacter
  \fi
  \sphinxDUC{00A0}{\nobreakspace}
  \sphinxDUC{2500}{\sphinxunichar{2500}}
  \sphinxDUC{2502}{\sphinxunichar{2502}}
  \sphinxDUC{2514}{\sphinxunichar{2514}}
  \sphinxDUC{251C}{\sphinxunichar{251C}}
  \sphinxDUC{2572}{\textbackslash}
\fi
\usepackage{cmap}
\usepackage[T1]{fontenc}
\usepackage{amsmath,amssymb,amstext}
\usepackage{babel}



\usepackage{tgtermes}
\usepackage{tgheros}
\renewcommand{\ttdefault}{txtt}



\usepackage[Bjarne]{fncychap}
\usepackage{sphinx}

\fvset{fontsize=auto}
\usepackage{geometry}


% Include hyperref last.
\usepackage{hyperref}
% Fix anchor placement for figures with captions.
\usepackage{hypcap}% it must be loaded after hyperref.
% Set up styles of URL: it should be placed after hyperref.
\urlstyle{same}

\addto\captionsenglish{\renewcommand{\contentsname}{Contents:}}

\usepackage{sphinxmessages}
\setcounter{tocdepth}{1}



\title{VITA}
\date{Sep 12, 2025}
\release{}
\author{Fernanda Pereira Guidotti Carneiro}
\newcommand{\sphinxlogo}{\vbox{}}
\renewcommand{\releasename}{}
\makeindex
\begin{document}

\ifdefined\shorthandoff
  \ifnum\catcode`\=\string=\active\shorthandoff{=}\fi
  \ifnum\catcode`\"=\active\shorthandoff{"}\fi
\fi

\pagestyle{empty}
\sphinxmaketitle
\pagestyle{plain}
\sphinxtableofcontents
\pagestyle{normal}
\phantomsection\label{\detokenize{index::doc}}


\sphinxAtStartPar
The \sphinxstylestrong{VITA Platform} (\sphinxstyleemphasis{Visionary Industrial Technology Architecture}) is designed as a modular and scalable environment to accelerate the development and deployment of AI solutions as services, with a strong focus on industrial applications and integration with the 6C Architecture.

\sphinxAtStartPar
VITA is implemented in \sphinxstylestrong{Python} and exposes its functionalities through a \sphinxstylestrong{REST API} using JSON for communication.

\sphinxAtStartPar
The documentation for the entire source code is provided below.


\chapter{Contents}
\label{\detokenize{index:contents}}
\sphinxstepscope


\section{INTRODUCTION}
\label{\detokenize{introduction:introduction}}\label{\detokenize{introduction::doc}}
\sphinxAtStartPar
The \sphinxstylestrong{VITA} Platform (Visionary Industrial Technology Architecture) was designed to enable the practical application of the 6C Architecture in industrial environments, offering a modular ecosystem for data analysis, forecasting, and process optimization.
Developed in Python, VITA adopts a flexible and scalable design, communicating through REST APIs in JSON format, ensuring interoperability with external systems and seamless integration with industrial pipelines.

\sphinxAtStartPar
The project follows a folder structure that facilitates development, maintenance, and evolution:
\begin{itemize}
\item {} 
\sphinxAtStartPar
\sphinxtitleref{algorithms} \textendash{} Contains forecasting and analytics algorithms, including LSTM, Reservoir Computing, hybrid models, and AutoML, each in its own folder for modularity and clarity.

\item {} 
\sphinxAtStartPar
\sphinxtitleref{common} \textendash{} Holds common modules and utility functions used across the application, such as path configuration, data handling, and logging.

\item {} 
\sphinxAtStartPar
\sphinxtitleref{dashboard} \textendash{} Implements VITA’s interactive interface, enabling visualization and analysis of forecasts and metrics.

\item {} 
\sphinxAtStartPar
\sphinxtitleref{datasets} \textendash{} Contains datasets used for model training, testing, and evaluation, including historical data and exogenous variables.

\item {} 
\sphinxAtStartPar
\sphinxtitleref{docs} \textendash{} Stores documentation generated with the Sphinx library (v4.0.1), describing the architecture, APIs, and processing flows.

\item {} 
\sphinxAtStartPar
\sphinxtitleref{log} \textendash{} Records execution logs, errors, and platform events for auditing and maintenance.

\item {} 
\sphinxAtStartPar
\sphinxtitleref{output\_forecasts} \textendash{} Holds forecasts generated by the models, organized by customer, algorithm, and scenario.

\item {} 
\sphinxAtStartPar
\sphinxtitleref{services} \textendash{} Contains the implementation of services and API endpoints that orchestrate the platform’s functionalities.

\item {} 
\sphinxAtStartPar
\sphinxtitleref{test} \textendash{} Contains scripts and test data for module and service validation.

\end{itemize}

\sphinxAtStartPar
Files in the project root include:
\begin{itemize}
\item {} 
\sphinxAtStartPar
\sphinxtitleref{app.py} \textendash{} Main application entry point.

\item {} 
\sphinxAtStartPar
\sphinxtitleref{config.py} \textendash{} General system configuration.

\item {} 
\sphinxAtStartPar
\sphinxtitleref{docker\sphinxhyphen{}compose.yml} and Dockerfile \textendash{} Container deployment configuration files.

\item {} 
\sphinxAtStartPar
\sphinxtitleref{requirements.txt} \textendash{} List of application dependencies.

\item {} 
\sphinxAtStartPar
\sphinxtitleref{setup.py} \textendash{} Installation and package configuration script.

\end{itemize}

\sphinxAtStartPar
This modular organization makes VITA a robust, scalable, and adaptable solution, addressing both research and development needs as well as real\sphinxhyphen{}world industrial deployments.

\sphinxAtStartPar
The general folder structure of the VITA Platform is shown below:

\noindent{\hspace*{\fill}\sphinxincludegraphics[scale=0.3]{{project-structure}.png}\hspace*{\fill}}

\sphinxstepscope


\section{ALGORITHMS}
\label{\detokenize{algorithms:algorithms}}\label{\detokenize{algorithms::doc}}
\sphinxAtStartPar
This section documents the internal code used for building and executing the Artificial Intelligence models within the VITA Platform.
All modules are organized inside the \sphinxtitleref{algorithms} directory in a modular way to facilitate maintenance, testing, and extensibility.

\sphinxAtStartPar
General Structure
\sphinxtitleref{core} \textendash{} Contains utility functions, data preprocessing routines, normalization methods, and shared components for all models.

\sphinxAtStartPar
\sphinxtitleref{hybrid\_model} \textendash{} Implements the hybrid model, which combines different forecasting approaches to improve accuracy.

\sphinxAtStartPar
\sphinxtitleref{lstm} \textendash{} Implementation of the Long Short\sphinxhyphen{}Term Memory model, used for multivariate time series forecasting.

\sphinxAtStartPar
\sphinxtitleref{pycaret} \textendash{} Integrates the PyCaret framework, enabling rapid experimentation with multiple algorithms.

\sphinxAtStartPar
\sphinxtitleref{rc} \textendash{} Implements Reservoir Computing, a recurrent neural network\sphinxhyphen{}based technique for time series forecasting.


\bigskip\hrule\bigskip



\subsection{\sphinxstylestrong{Module Documentation}}
\label{\detokenize{algorithms:module-documentation}}
\sphinxAtStartPar
\sphinxtitleref{Core}

\sphinxAtStartPar
\sphinxtitleref{Hybrid Model}

\sphinxAtStartPar
\sphinxtitleref{LSTM}

\sphinxAtStartPar
\sphinxtitleref{PyCaret}

\sphinxAtStartPar
\sphinxtitleref{Reservoir Computing (RC)}

\sphinxstepscope


\section{ARCHITECTURE 6C}
\label{\detokenize{architecture6c:architecture-6c}}\label{\detokenize{architecture6c::doc}}
\sphinxAtStartPar
The \sphinxstylestrong{6C Architecture} provides a conceptual framework for building intelligent and connected
industrial systems. It defines six layers that structure the flow from raw data acquisition
to advanced decision\sphinxhyphen{}making, ensuring adaptability, traceability, and continuous improvement.

\sphinxAtStartPar
VITA adopts this architecture as a guiding principle, aligning its modular components with each of the six layers.


\subsection{Conceptual Representation}
\label{\detokenize{architecture6c:conceptual-representation}}
\sphinxAtStartPar
The classical representation of the 6C model is shown below. It highlights the six layers:
Connection, Conversion, Cyber, Cognition, Configuration, and Consciousness.

\noindent{\hspace*{\fill}\sphinxincludegraphics[scale=0.08]{{plattform}.png}\hspace*{\fill}}


\subsection{Applied Industrial View}
\label{\detokenize{architecture6c:applied-industrial-view}}
\sphinxAtStartPar
Beyond the conceptual pyramid, the 6C Architecture can also be visualized as an operational
pipeline, showing how industrial data flows across layers. This representation emphasizes
data ingestion, transformation, modeling, decision support, and continuous self\sphinxhyphen{}adaptation.

\noindent{\hspace*{\fill}\sphinxincludegraphics[scale=0.5]{{arquitetura6c}.png}\hspace*{\fill}}


\subsection{Mapping to VITA}
\label{\detokenize{architecture6c:mapping-to-vita}}
\sphinxAtStartPar
Each layer of the 6C Architecture corresponds to specific modules within the VITA Platform:
\begin{itemize}
\item {} 
\sphinxAtStartPar
\sphinxstylestrong{Connection} \(\rightarrow\) \sphinxtitleref{datasets}, \sphinxtitleref{services} (data ingestion and standardized access)

\item {} 
\sphinxAtStartPar
\sphinxstylestrong{Conversion} \(\rightarrow\) \sphinxtitleref{common}, preprocessing routines (normalization and preparation)

\item {} 
\sphinxAtStartPar
\sphinxstylestrong{Cyber} \(\rightarrow\) \sphinxtitleref{algorithms} (LSTM, Reservoir Computing, Hybrid, PyCaret)

\item {} 
\sphinxAtStartPar
\sphinxstylestrong{Cognition} \(\rightarrow\) \sphinxtitleref{output\_forecasts}, performance metrics (knowledge discovery and insights)

\item {} 
\sphinxAtStartPar
\sphinxstylestrong{Configuration} \(\rightarrow\) \sphinxtitleref{dashboard} (interactive visualization, scenario setup)

\item {} 
\sphinxAtStartPar
\sphinxstylestrong{Consciousness} \(\rightarrow\) \sphinxtitleref{log}, monitoring (continuous learning, evaluation, auditing, self\sphinxhyphen{}adaptation)

\end{itemize}

\sphinxAtStartPar
This dual perspective (conceptual and applied) shows how VITA is both academically grounded
and practically designed, bridging the gap between theory and industrial deployment.

\sphinxstepscope


\section{DATA FLOW, CRISP\sphinxhyphen{}DM \& MLOps}
\label{\detokenize{dataflow:data-flow-crisp-dm-mlops}}\label{\detokenize{dataflow::doc}}

\subsection{Overview}
\label{\detokenize{dataflow:overview}}
\sphinxAtStartPar
The \sphinxstylestrong{VITA Platform} (\sphinxstyleemphasis{Visionary Industrial Technology Architecture}) integrates
three complementary perspectives:
\begin{itemize}
\item {} 
\sphinxAtStartPar
\sphinxstylestrong{CRISP\sphinxhyphen{}DM} \textendash{} methodological workflow for analytics and forecasting.

\item {} 
\sphinxAtStartPar
\sphinxstylestrong{6C Architecture} \textendash{} conceptual industrial intelligence framework.

\item {} 
\sphinxAtStartPar
\sphinxstylestrong{MLOps} \textendash{} operational lifecycle for machine learning in production.

\end{itemize}

\sphinxAtStartPar
Together, they ensure that VITA is not only research\sphinxhyphen{}driven, but also scalable,
auditable, and production\sphinxhyphen{}ready.


\subsection{Integrated Mapping}
\label{\detokenize{dataflow:integrated-mapping}}

\begin{savenotes}\sphinxattablestart
\sphinxthistablewithglobalstyle
\centering
\begin{tabulary}{\linewidth}[t]{TTTT}
\sphinxtoprule
\sphinxstyletheadfamily 
\sphinxAtStartPar
CRISP\sphinxhyphen{}DM Phase
&\sphinxstyletheadfamily 
\sphinxAtStartPar
6C Layer
&\sphinxstyletheadfamily 
\sphinxAtStartPar
MLOps Perspective
&\sphinxstyletheadfamily 
\sphinxAtStartPar
VITA Modules / Artifacts
\\
\sphinxmidrule
\sphinxtableatstartofbodyhook
\sphinxAtStartPar
Business Understanding
&
\sphinxAtStartPar
Configuration
&
\sphinxAtStartPar
Experiment tracking
&
\sphinxAtStartPar
\sphinxtitleref{dashboard}, docs
\\
\sphinxhline
\sphinxAtStartPar
Data Understanding
&
\sphinxAtStartPar
Connection
&
\sphinxAtStartPar
DataOps (ingestion)
&
\sphinxAtStartPar
\sphinxtitleref{services}, \sphinxtitleref{datasets}
\\
\sphinxhline
\sphinxAtStartPar
Data Preparation
&
\sphinxAtStartPar
Conversion
&
\sphinxAtStartPar
DataOps (preprocessing)
&
\sphinxAtStartPar
\sphinxtitleref{common}, \sphinxtitleref{algorithms/core}
\\
\sphinxhline
\sphinxAtStartPar
Modeling
&
\sphinxAtStartPar
Cyber
&
\sphinxAtStartPar
ModelOps (training)
&
\sphinxAtStartPar
\sphinxtitleref{algorithms/lstm}, \sphinxtitleref{algorithms/rc},
\sphinxtitleref{algorithms/hybrid\_model}, \sphinxtitleref{algorithms/pycaret}
\\
\sphinxhline
\sphinxAtStartPar
Evaluation
&
\sphinxAtStartPar
Cognition
&
\sphinxAtStartPar
Model validation
&
\sphinxAtStartPar
\sphinxtitleref{output\_forecasts}, metrics, reports
\\
\sphinxhline
\sphinxAtStartPar
Deployment
&
\sphinxAtStartPar
Configuration
&
\sphinxAtStartPar
CI/CD pipelines
&
\sphinxAtStartPar
\sphinxtitleref{services} (inference APIs), \sphinxtitleref{dashboard} (visual rollout)
\\
\sphinxhline
\sphinxAtStartPar
Monitoring \& Feedback
&
\sphinxAtStartPar
Consciousness
&
\sphinxAtStartPar
Continuous learning
&
\sphinxAtStartPar
\sphinxtitleref{log}, monitoring hooks, retraining triggers
\\
\sphinxbottomrule
\end{tabulary}
\sphinxtableafterendhook\par
\sphinxattableend\end{savenotes}


\subsection{End\sphinxhyphen{}to\sphinxhyphen{}End Flow}
\label{\detokenize{dataflow:end-to-end-flow}}\begin{enumerate}
\sphinxsetlistlabels{\arabic}{enumi}{enumii}{}{.}%
\item {} 
\sphinxAtStartPar
\sphinxstylestrong{Ingest \& Profile} data (\sphinxtitleref{services} \(\rightarrow\) \sphinxtitleref{datasets}).

\item {} 
\sphinxAtStartPar
\sphinxstylestrong{Prepare \& Convert} with \sphinxtitleref{common} and \sphinxtitleref{algorithms/core}.

\item {} 
\sphinxAtStartPar
\sphinxstylestrong{Model} with \sphinxtitleref{algorithms} (LSTM, RC, Hybrid, PyCaret).

\item {} 
\sphinxAtStartPar
\sphinxstylestrong{Evaluate} results, save to \sphinxtitleref{output\_forecasts}.

\item {} 
\sphinxAtStartPar
\sphinxstylestrong{Deploy} via \sphinxtitleref{services} APIs and \sphinxtitleref{dashboard}.

\item {} 
\sphinxAtStartPar
\sphinxstylestrong{Monitor} through \sphinxtitleref{log/}, detect drift, close the loop.

\item {} 
\sphinxAtStartPar
\sphinxstylestrong{Retrain automatically} (MLOps) when monitoring triggers thresholds.

\end{enumerate}


\subsection{Implementation Notes}
\label{\detokenize{dataflow:implementation-notes}}\begin{itemize}
\item {} 
\sphinxAtStartPar
\sphinxstylestrong{DataOps} in VITA: ingestion (\sphinxtitleref{services}), datasets, preprocessing (\sphinxtitleref{common}).

\item {} 
\sphinxAtStartPar
\sphinxstylestrong{ModelOps} in VITA: modular algorithms, experiment tracking, forecasts.

\item {} 
\sphinxAtStartPar
\sphinxstylestrong{CI/CD}: Docker\sphinxhyphen{}based deployments (\sphinxtitleref{docker\sphinxhyphen{}compose.yml}, \sphinxtitleref{Dockerfile}).

\item {} 
\sphinxAtStartPar
\sphinxstylestrong{Monitoring}: \sphinxtitleref{log/} + retraining pipelines close the loop for industrial AI.

\end{itemize}

\sphinxstepscope


\section{GLOSSARY}
\label{\detokenize{glossary:glossary}}\label{\detokenize{glossary::doc}}\begin{description}
\sphinxlineitem{6C Architecture\index{6C Architecture@\spxentry{6C Architecture}|spxpagem}\phantomsection\label{\detokenize{glossary:term-6C-Architecture}}}
\sphinxAtStartPar
Conceptual industrial intelligence framework composed of six layers:
Connection, Conversion, Cyber, Cognition, Configuration, and Consciousness.

\sphinxlineitem{API (Application Programming Interface)\index{API (Application Programming Interface)@\spxentry{API}\spxextra{Application Programming Interface}|spxpagem}\phantomsection\label{\detokenize{glossary:term-API-Application-Programming-Interface}}}
\sphinxAtStartPar
A set of rules and protocols allowing software components to communicate
with each other. In VITA, APIs are implemented with FastAPI and expose
training and forecasting endpoints.

\sphinxlineitem{CRISP\sphinxhyphen{}DM\index{CRISP\sphinxhyphen{}DM@\spxentry{CRISP\sphinxhyphen{}DM}|spxpagem}\phantomsection\label{\detokenize{glossary:term-CRISP-DM}}}
\sphinxAtStartPar
Cross\sphinxhyphen{}Industry Standard Process for Data Mining. A methodology that defines
phases for analytics projects: Business Understanding, Data Understanding,
Data Preparation, Modeling, Evaluation, and Deployment.

\sphinxlineitem{Dashboard\index{Dashboard@\spxentry{Dashboard}|spxpagem}\phantomsection\label{\detokenize{glossary:term-Dashboard}}}
\sphinxAtStartPar
Interactive web interface (Streamlit\sphinxhyphen{}based) of VITA that allows users to
explore forecasts, configure scenarios, and visualize performance metrics.

\sphinxlineitem{DataOps\index{DataOps@\spxentry{DataOps}|spxpagem}\phantomsection\label{\detokenize{glossary:term-DataOps}}}
\sphinxAtStartPar
Practices that focus on improving the communication, integration, and
automation of data flows across the data lifecycle.

\sphinxlineitem{Dataset\index{Dataset@\spxentry{Dataset}|spxpagem}\phantomsection\label{\detokenize{glossary:term-Dataset}}}
\sphinxAtStartPar
Structured collection of historical records and exogenous variables used
for training and testing forecasting models.

\sphinxlineitem{Exogenous Variables\index{Exogenous Variables@\spxentry{Exogenous Variables}|spxpagem}\phantomsection\label{\detokenize{glossary:term-Exogenous-Variables}}}
\sphinxAtStartPar
External factors that influence the target variable. In VITA examples,
economic indicators (e.g., SELIC, IPCA) or production orders.

\sphinxlineitem{Forecast\index{Forecast@\spxentry{Forecast}|spxpagem}\phantomsection\label{\detokenize{glossary:term-Forecast}}}
\sphinxAtStartPar
Predicted value generated by a model. VITA forecasts are saved in
\sphinxcode{\sphinxupquote{output\_forecasts/}}.

\sphinxlineitem{Forecast Types\index{Forecast Types@\spxentry{Forecast Types}|spxpagem}\phantomsection\label{\detokenize{glossary:term-Forecast-Types}}}\begin{itemize}
\item {} 
\sphinxAtStartPar
\sphinxstylestrong{Validation} \textendash{} Backtest forecasts on the last 4 months of historical data
(holdout set), used to evaluate model accuracy.

\item {} 
\sphinxAtStartPar
\sphinxstylestrong{Validation+OS} \textendash{} Validation forecasts with \sphinxstyleemphasis{outliers substituted}
before evaluation.

\item {} 
\sphinxAtStartPar
\sphinxstylestrong{Production} \textendash{} Forecasts for the next 4 months after training on the
full historical dataset.

\item {} 
\sphinxAtStartPar
\sphinxstylestrong{Production+OS} \textendash{} Production forecasts with \sphinxstyleemphasis{outliers substituted}
before forecasting.

\item {} 
\sphinxAtStartPar
\sphinxstylestrong{Test} \textendash{} Dataset of the last 4 historical months with real demand
values and forecasts for evaluation.

\item {} 
\sphinxAtStartPar
\sphinxstylestrong{Future} \textendash{} Forecasts for the next 4 months without real demand
(prediction horizon).

\item {} 
\sphinxAtStartPar
\sphinxstylestrong{Variation} \textendash{} Different feature set combinations in \sphinxcode{\sphinxupquote{model\_input}}
(e.g., \sphinxcode{\sphinxupquote{cd}}, \sphinxcode{\sphinxupquote{cd\_po}}, \sphinxcode{\sphinxupquote{cd\_po\_selic\_ipca}}).

\end{itemize}

\sphinxlineitem{Inputs (Features)\index{Inputs (Features)@\spxentry{Inputs}\spxextra{Features}|spxpagem}\phantomsection\label{\detokenize{glossary:term-Inputs-Features}}}\begin{itemize}
\item {} 
\sphinxAtStartPar
\sphinxstylestrong{cd} \textendash{} \sphinxstyleemphasis{Customer Demand} (historical demand, target variable).

\item {} 
\sphinxAtStartPar
\sphinxstylestrong{po} \textendash{} \sphinxstyleemphasis{Production Orders}.

\item {} 
\sphinxAtStartPar
\sphinxstylestrong{ipca} \textendash{} \sphinxstyleemphasis{Brazilian inflation index (IPCA)}.

\item {} 
\sphinxAtStartPar
\sphinxstylestrong{selic} \textendash{} \sphinxstyleemphasis{Brazilian interest rate (SELIC)}.

\item {} 
\sphinxAtStartPar
\sphinxstylestrong{model\_input} \textendash{} Feature set used to train the model (e.g., \sphinxcode{\sphinxupquote{cd\_po\_ipca}}).

\item {} 
\sphinxAtStartPar
\sphinxstylestrong{forecast\_ai} \textendash{} Predicted values generated by the model.

\item {} 
\sphinxAtStartPar
\sphinxstylestrong{customer\_demand} \textendash{} Real observed demand values.

\item {} 
\sphinxAtStartPar
\sphinxstylestrong{forecast\_deviation\_abs} \textendash{} Absolute forecast error, usually expressed in \%.

\end{itemize}

\sphinxlineitem{LSTM (Long Short\sphinxhyphen{}Term Memory)\index{LSTM (Long Short\sphinxhyphen{}Term Memory)@\spxentry{LSTM}\spxextra{Long Short\sphinxhyphen{}Term Memory}|spxpagem}\phantomsection\label{\detokenize{glossary:term-LSTM-Long-Short-Term-Memory}}}
\sphinxAtStartPar
A type of recurrent neural network designed to capture long\sphinxhyphen{}term dependencies
in sequential data. Used for multivariate time series forecasting.

\sphinxlineitem{MLOps\index{MLOps@\spxentry{MLOps}|spxpagem}\phantomsection\label{\detokenize{glossary:term-MLOps}}}
\sphinxAtStartPar
Set of practices for deploying and maintaining machine learning models in
production reliably and efficiently. Includes CI/CD, monitoring, retraining,
and model versioning.

\sphinxlineitem{Module\index{Module@\spxentry{Module}|spxpagem}\phantomsection\label{\detokenize{glossary:term-Module}}}
\sphinxAtStartPar
Self\sphinxhyphen{}contained part of the platform codebase, usually corresponding to a
functional unit (e.g., \sphinxcode{\sphinxupquote{algorithms}}, \sphinxcode{\sphinxupquote{services}}, \sphinxcode{\sphinxupquote{dashboard}}).

\sphinxlineitem{Reservoir Computing (RC)\index{Reservoir Computing (RC)@\spxentry{Reservoir Computing}\spxextra{RC}|spxpagem}\phantomsection\label{\detokenize{glossary:term-Reservoir-Computing-RC}}}
\sphinxAtStartPar
A recurrent neural network approach that leverages a randomly connected
dynamic reservoir for efficient time series forecasting.

\sphinxlineitem{Scenario\index{Scenario@\spxentry{Scenario}|spxpagem}\phantomsection\label{\detokenize{glossary:term-Scenario}}}
\sphinxAtStartPar
A configuration that defines the target variable and exogenous inputs used
in a forecasting task. Example: \sphinxcode{\sphinxupquote{cd\_po\_selic\_ipca}} indicates
\sphinxstyleemphasis{customer\_demand} as target with \sphinxstyleemphasis{production\_orders}, \sphinxstyleemphasis{SELIC}, and \sphinxstyleemphasis{IPCA}
as exogenous variables.

\sphinxlineitem{Service\index{Service@\spxentry{Service}|spxpagem}\phantomsection\label{\detokenize{glossary:term-Service}}}
\sphinxAtStartPar
A microservice implemented with FastAPI that provides functionalities such
as data upload, training, and prediction.

\sphinxlineitem{VITA Platform\index{VITA Platform@\spxentry{VITA Platform}|spxpagem}\phantomsection\label{\detokenize{glossary:term-VITA-Platform}}}
\sphinxAtStartPar
\sphinxstylestrong{Visionary Industrial Technology Architecture}. A modular ecosystem for
data analysis, forecasting, and process optimization in industrial contexts.

\end{description}


\chapter{Indices and Tables}
\label{\detokenize{index:indices-and-tables}}\begin{itemize}
\item {} 
\sphinxAtStartPar
\DUrole{xref}{\DUrole{std}{\DUrole{std-ref}{genindex}}}

\item {} 
\sphinxAtStartPar
\DUrole{xref}{\DUrole{std}{\DUrole{std-ref}{modindex}}}

\item {} 
\sphinxAtStartPar
\DUrole{xref}{\DUrole{std}{\DUrole{std-ref}{search}}}

\end{itemize}



\renewcommand{\indexname}{Index}
\printindex
\end{document}